\documentclass[../../../include/open-logic-section]{subfiles}

\begin{document}

\olfileid{sth}{choice}{tarskiscott}
\olsection{تارسکی–سکاٹ دی ترکیب}

\olref[cardinals][cardsasords]{defcardinalasordinal} وچ اسیں کارڈینل عدداں دی تعریف ترتیبی عدداں وجوں دتی سی۔ ایس کم لئی اسیں حُسنِ ترتیب دا مسلّمہ منّیا سی۔ ایس دی وجہ بس ایہ سی کہ شعبے وچ ایہو «چلدا معیار» اے۔

سو حُسنِ ترتیب نال جڑے فلسفیانہ مسئلیاں دی گل توں پہلاں، ایہ صاف ہونا ضروری اے کہ اسیں چلدا معیار چھڈ وی \emph{سکدے آں} تے حُسنِ ترتیب منّے \emph{بناں} کارڈینل عدداں دا نظریہ بنا سکدے آں۔ اسیں ہن وی $\cardeq{A}{B}$، $\cardle{A}{B}$ تے $\cardless{A}{B}$ دیاں اوہ تعریفاں ورت سکدے آں جیہڑیاں \olref[sfr][siz][]{chap} وچ دتیاں سن۔ سانوں بس \emph{کارڈینل عدد} دا نواں تصور چاہیدا اے۔

اک سادہ خیال ایہ ہو سکدا اے کہ $A$ دی کارڈینیلٹی ایویں تعریف کرن دی کوشش کریے:
\[
	\Setabs{x}{\cardeq{A}{x}}.
\]
تسی ایہدا مقابلہ فریگے دی $\# x Fx$ والی تعریف نال کر سکدے او، جیہدا خاکہ \olref[cardinals][hp]{sec} دے عین آخر وچ دتا سی۔ تے جیہڑیاں وجوہاں دا اوتھے اشارہ دتا سی، اوہناں کرکے ایہ تعریف نہیں چل دی۔ ہر اک رکنی سیٹ $\{\emptyset\}$ نال ہم تعداد اے۔ پر درجہ وار سلسلے دے ہر جانشین مرحلے اُتے نویں اک رکنی سیٹ بندے نیں (بس اوہ اک رکنی سیٹ ویکھو جیہدا اکّو رکن پچھلا مرحلہ اے)۔ سو $\Setabs{x}{\cardeq{A}{x}}$ موجود نہیں، کیوں جے ایہدا کوئی درجہ نہیں ہو سکدا۔

ایس اوکڑ توں لنگھن لئی اسیں تارسکی تے سکاٹ دی اک ترکیب ورتدے آں۔\footnote{چیتے رہے کہ سارے فارمولیاں وچ پیرامیٹر ہو سکدے نیں (جے صاف ہور گل نہ آکھی ہووے)۔}

\begin{defn}[تارسکی–سکاٹ]
کسے وی فارمولے $\phi(x)$ لئی، $[ x : \phi(x)] $ اوہناں سارے $x$ دا سیٹ ہووے جیہناں دا درجہ سبھ توں گھٹ ممکن درجہ اے تے جیہناں لئی $\phi(x)$ اے (یا $\emptyset$، جے کوئی $\phi$ پوری کرن والی شے نہیں)۔
\end{defn}

سانوں جانچنا چاہیدا اے کہ ایہ تعریف جائز اے۔ $\ZF$ وچ کم کردیاں، \olref[spine][foundation]{zfentailsregularity} ضمانت دیندا اے کہ $\setrank{x}$ ہر $x$ لئی موجود اے۔ ہُن جے کوئی شےواں $\phi$ اُتے پوریاں اُتردیاں نیں، تاں $\alpha$ اوہ سبھ توں گھٹ درجہ ہووے جیہدے لئی $(\exists x\subseteq V_\alpha)\phi(x)$، یعنی $(\forall \beta
\in \alpha)(\forall x \subseteq V_\beta)\lnot \phi(x)$۔ فیر وکھ کرن دے اصول نال $[x : \phi(x)]$ دی تعریف $\Setabs{x \in
V_{\alpha+1}}{\phi(x)}$ وجوں دے سکدے آں۔

تارسکی–سکاٹ دی ترکیب دا جواز دین پچھوں، ہُن اسیں ایہنوں کارڈینیلٹی دا اک تصور تعریف کرن لئی ورت سکدے آں۔

\begin{defn}
$A$ دی \textsc{ts}-کارڈینیلٹی $\text{tsc}(A) = [x :
\cardeq{A}{x}]$ اے۔
\end{defn}

اک \textsc{ts}-کارڈینل عدد دی تعریف وچ حُسنِ ترتیب دا مسلّمہ نہیں ورتیا جاندا۔ پر ایس مسلّمے توں بناں وی اسیں وکھا سکدے آں کہ \emph{\textsc{ts}-کارڈینل عدد}، \olref[cardinals][cardsasords]{defcardinalasordinal} وچ دتی تعریف والے \emph{کارڈینل عدداں} وانگ کم کردے نیں۔ مثلاً جے اسیں \olref[cardinals][cardsasords]{lem:CardinalsBehaveRight} تے \olref[card-arithmetic][opps]{lem:SizePowerset2Exp} نوں \textsc{ts}-کارڈینل عدداں دے حساب نال فیر لکھیے، تاں ثبوت $\ZF$ وچ ٹھیک چلدے نیں، حُسنِ ترتیب دا مسلّمہ منّے بناں۔

ایہدی گل کردیاں دھیان دیو کہ اسیں تارسکی–سکاٹ دی ترکیب ورت کے ترتیبی عدداں دا نظریہ وی بنا سکدے آں۔ جے $\tuple{A, <}$ اک حُسنِ ترتیب اے، تاں $\text{tso}(A, <) = [\tuple{X,
R} : \ordeq{\tuple{A, <}}{\tuple{X, R}}]$ ہووے۔ کارڈینل تے ترتیبی عدداں دے ایس بیان بارے ہور لئی \citet[chs.~9--12]{Potter2004} ویکھو۔

\end{document}
